% ============================================================
%  Chapter 4 — Results
% ============================================================
\chapter{Results}
\label{ch:results}

This chapter reports the outcomes of the analysis in the order the methodology
introduced them: the exploratory structure of the feature data, the comparative
performance of the four classifiers, the factors the best model relied on, the
resulting groundwater potential map, and finally the spatial-realism validation
that tests the map against independent evidence.

\section{Exploratory Data Analysis}
\label{sec:eda}

Before any model was trained, the relationship between each feature and the
groundwater potential label was examined across all \npixels{} valid pixels. The
nine features contained no missing values within the study-area mask and each
spanned the full normalized range, so no pixels were lost to incomplete data.

The Pearson correlations (Figure~\ref{fig:correlation}) are informative in their
own right. Soil texture stands out, correlating with the potential label at
$r = 0.836$---far stronger than any other feature---followed by cumulative
rainfall at $r = 0.514$. Land surface temperature is the strongest of the
negatively associated features ($r = -0.426$), consistent with its role as a
proxy for evaporative loss. NDVI ($+0.240$), elevation ($+0.234$), NDWI
($-0.222$), and TWI ($+0.185$) form a middle group, while aspect ($+0.060$) and
slope ($+0.051$) are effectively uncorrelated. The negligible slope correlation
is itself a regional finding rather than a defect: across a comparatively
subdued plateau, slope simply does not discriminate groundwater zones the way it
would in mountainous terrain.

\begin{figure}[H]
  \centering
  \includegraphics[width=0.8\textwidth]{correlation_matrix}
  \caption{Pearson correlation matrix of the nine input features and the
           groundwater potential label. Soil texture ($r=0.836$) and rainfall
           ($r=0.514$) are the dominant positively correlated predictors.}
  \label{fig:correlation}
\end{figure}

The dominance of soil and rainfall is consistent with the hydrogeology discussed
in Chapter~\ref{ch:methodology} and with the wider literature: permeability
controls whether water infiltrates, and rainfall supplies the water to begin
with. It also echoes Khan et al. \cite{khan2023rainfallrunoff}, who found
rainfall to be the principal driver of the hydrological response in these same
basins. Figure~\ref{fig:feature_target} summarizes the per-feature association
with the label in ranked form.

\begin{figure}[H]
  \centering
  \includegraphics[width=0.75\textwidth]{feature_target_correlation}
  \caption{Features ranked by absolute correlation with the groundwater potential
           label, showing the clear lead of soil texture and rainfall.}
  \label{fig:feature_target}
\end{figure}

By construction, the weighted-overlay labelling produced a class distribution of
40\% Low, 40\% Medium, and 20\% High (Figure~\ref{fig:class_dist}). The balanced
training sample drawn from these labels removed this imbalance for model fitting,
as described in Section~\ref{sec:training_data}.

\begin{figure}[H]
  \centering
  \includegraphics[width=0.7\textwidth]{class_distribution_pub}
  \caption{Distribution of the three groundwater potential classes in the
           weighted-overlay reference labels across the full plateau.}
  \label{fig:class_dist}
\end{figure}

\section{Model Performance}
\label{sec:model_performance}

All four classifiers were trained on the 360,000-pixel training portion and
evaluated on the held-out 90,000-pixel test set.
Table~\ref{tab:model_results} reports the headline metrics.

\begin{table}[H]
\centering
\caption{Performance of the four classifiers on the 90,000-pixel test set. Best
         value in each column in bold.}
\label{tab:model_results}
\begin{tabular}{lcccc}
\toprule
Model & Accuracy (\%) & Macro $F_1$ & Cohen's $\kappa$ & ROC-AUC \\
\midrule
Random Forest      & 98.32 & 0.9832 & 0.9749 & 0.9995 \\
\textbf{XGBoost}   & \textbf{99.12} & \textbf{0.9912} & \textbf{0.9868} & \textbf{0.9998} \\
LightGBM           & 99.11 & 0.9911 & 0.9867 & 0.9998 \\
Voting Ensemble    & 99.09 & 0.9909 & 0.9863 & 0.9998 \\
\bottomrule
\end{tabular}
\end{table}

XGBoost gave the best results on every metric, reaching 99.12\% accuracy, a macro
$F_1$ of 0.9912, Cohen's $\kappa$ of 0.9868, and a one-vs-rest ROC-AUC of 0.9998,
and it was selected as the model for full-plateau prediction. The margins between
the top three models are slim---LightGBM trails by barely a hundredth of a
percentage point, and the voting ensemble sits just behind both---while random
forest is a clear, if still strong, step below at 98.32\%. This ordering, with
XGBoost ahead of random forest and gradient boosting clustered at the top, mirrors
the ranking Sahin \cite{sahin2020ensembletree} reported for the same model
families in landslide susceptibility mapping, which lends some external
credibility to the comparison even though the absolute numbers here are higher.
Figure~\ref{fig:model_comparison} shows the comparison visually.

\begin{figure}[H]
  \centering
  \includegraphics[width=0.85\textwidth]{model_comparison}
  \caption{Comparison of the four classifiers across accuracy, macro $F_1$,
           Cohen's $\kappa$, and ROC-AUC on the test set.}
  \label{fig:model_comparison}
\end{figure}

The absolute accuracies---near 99\%---are far above the 0.73--0.86 AUC range
typical of knowledge-driven AHP and frequency-ratio GWPM studies
\cite{das2018pravara,pande2021ahpmif}. This gap should be read carefully rather
than celebrated uncritically: it reflects in large part the fact that the models
are learning to reproduce a deterministic weighted overlay built from the same
features, a point developed in Section~\ref{sec:agreement} and again in
Chapter~\ref{ch:discussion}. What the metrics establish is that the learning
problem was solved cleanly; whether the resulting map is hydrogeologically
trustworthy is a separate question, addressed by the validation in
Section~\ref{sec:validation_results}.

\subsection{Per-Class Performance}
\label{sec:per_class}

Because a high overall accuracy can still hide weakness on a minority class, the
per-class behaviour of XGBoost was examined (Table~\ref{tab:per_class}). All
three classes are predicted with precision and recall above 98\%. The High class,
the one of greatest practical interest, is recovered with the highest recall of
the three (99.61\%), meaning the model rarely misses genuinely high-potential
ground in the test set; the slightly lower precision and recall on the Medium
class reflect the expected difficulty of the middle category, whose boundaries
with Low and High are the least sharp.

\begin{table}[H]
\centering
\caption{Per-class precision, recall, and $F_1$ for the XGBoost model on the
         test set.}
\label{tab:per_class}
\begin{tabular}{lccc}
\toprule
Class & Precision (\%) & Recall (\%) & $F_1$ (\%) \\
\midrule
Low    & 99.32 & 99.43 & 99.38 \\
Medium & 99.04 & 98.32 & 98.68 \\
High   & 98.99 & 99.61 & 99.31 \\
\bottomrule
\end{tabular}
\end{table}

\subsection{Class Separability}
\label{sec:roc}

The one-vs-rest ROC curves (Figure~\ref{fig:roc}) sit almost on the top-left
corner for every class and model, corresponding to the AUC values near 0.9998 in
Table~\ref{tab:model_results}. The precision--recall curves
(Figure~\ref{fig:pr}) tell the same story from the complementary angle. While
such near-perfect separability would be implausible for a map validated against
field data, it is exactly what one expects when the target is a deterministic
function of the inputs; the curves confirm that the models have recovered that
function with very little residual confusion.

\begin{figure}[H]
  \centering
  \includegraphics[width=0.78\textwidth]{roc_curves}
  \caption{One-vs-rest ROC curves for the four classifiers on the test set. All
           curves approach the ideal upper-left corner (AUC $\approx 0.9998$).}
  \label{fig:roc}
\end{figure}

\begin{figure}[H]
  \centering
  \includegraphics[width=0.78\textwidth]{pr_curves}
  \caption{Precision--recall curves for the three potential classes, confirming
           high precision is maintained across the full recall range.}
  \label{fig:pr}
\end{figure}

\subsection{Cross-Validation}
\label{sec:cv}

To confirm that the test-set performance was not an artefact of one particular
split, each model was also assessed under five-fold cross-validation on the
training data (Table~\ref{tab:cv}). The fold-to-fold variation is very small for
all three tree ensembles---standard deviations of a few ten-thousandths in macro
$F_1$---indicating stable, reproducible performance rather than a fortunate
partition. LightGBM showed the tightest spread, XGBoost the best test-set
metrics; in practice the two are interchangeable on this data, and either would
serve as the production model.

\begin{table}[H]
\centering
\caption{Five-fold cross-validated macro $F_1$ (mean $\pm$ standard deviation)
         on the training set.}
\label{tab:cv}
\begin{tabular}{lc}
\toprule
Model & Macro $F_1$ (mean $\pm$ SD) \\
\midrule
Random Forest & $0.9839 \pm 0.0005$ \\
XGBoost       & $0.9896 \pm 0.0004$ \\
LightGBM      & $0.9911 \pm 0.0002$ \\
\bottomrule
\end{tabular}
\end{table}

\begin{figure}[H]
  \centering
  \includegraphics[width=0.8\textwidth]{confusion_matrices}
  \caption{Confusion matrices for the four classifiers on the test set. Almost
           all mass lies on the diagonal; the few off-diagonal errors occur
           between adjacent classes (Low--Medium and Medium--High).}
  \label{fig:confusion}
\end{figure}

The confusion matrices (Figure~\ref{fig:confusion}) make the error structure
explicit: what little misclassification occurs is confined to neighbouring
classes, with essentially no Low pixels mistaken for High or vice versa. For an
ordinal target this is the most benign error pattern possible.

\section{Feature Importance}
\label{sec:feature_importance}

The features the models relied on agree closely with the exploratory correlations.
For both random forest and XGBoost, soil texture was overwhelmingly the most
important predictor, with an XGBoost importance of 0.831; rainfall ranked second
(random forest importance 0.198, LightGBM 0.236). The remaining features carried
modest weight, and LightGBM---consistent with its leaf-wise growth---spread
importance somewhat more evenly than the other two models.
Figure~\ref{fig:feature_importance} shows the comparison.

\begin{figure}[H]
  \centering
  \includegraphics[width=\textwidth]{feature_importance}
  \caption{Feature importance across random forest, XGBoost, and LightGBM. Soil
           texture dominates, with rainfall a consistent second.}
  \label{fig:feature_importance}
\end{figure}

That a learned, nonlinear model independently arrives at the same two dominant
factors as the correlation analysis---soil and rainfall---is reassuring. It also
aligns with the reference study \cite{waqas2022chakwal}, in which geology and
land use (the closest analogues to soil permeability) carried the largest
weights, and it confirms that the model's behaviour is hydrogeologically
sensible rather than being driven by spurious associations.

\section{Groundwater Potential Map}
\label{sec:prediction_map}

Applying the XGBoost model to all \npixels{} valid pixels produced the
full-resolution groundwater potential map in Figure~\ref{fig:gw_map}. The spatial
distribution of the three classes is given in Table~\ref{tab:pixel_dist}.

\begin{table}[H]
\centering
\caption{Predicted groundwater potential distribution across the \studyarea{}.}
\label{tab:pixel_dist}
\begin{tabular}{lcc}
\toprule
Zone & Pixel count & Coverage (\%) \\
\midrule
Low potential    & 14{,}806{,}035 & 40.0 \\
Medium potential & 14{,}806{,}036 & 40.0 \\
High potential   &  7{,}403{,}020 & 20.0 \\
\midrule
\textbf{Total}   & \textbf{\npixels{}} & \textbf{100.0} \\
\bottomrule
\end{tabular}
\end{table}

\begin{figure}[H]
  \centering
  \includegraphics[width=\textwidth]{Fig1_GW_Prediction_Map}
  \caption{Groundwater potential map of the \studyarea{} predicted by XGBoost at
           30~m resolution. High-potential ground (green) concentrates along the
           northern alluvial tracts and major river valleys; low-potential ground
           (red) dominates the southern hard-rock terrain.}
  \label{fig:gw_map}
\end{figure}

The map places high potential where hydrogeology would lead one to expect it:
along the alluvial corridors of the Soan and the river valleys, and across the
wetter northern districts, while the folded and crystalline terrain of the south
and the Salt Range is predominantly Low. The continuous High-probability surface
(Figure~\ref{fig:gw_prob}) conveys the same pattern with its gradations intact,
and is the layer used for the validation in Section~\ref{sec:validation_results}.

\begin{figure}[H]
  \centering
  \includegraphics[width=\textwidth]{Fig2_GW_Probability_Map}
  \caption{Continuous probability of the High potential class from the XGBoost
           model, preserving the gradational detail lost in the hard
           classification.}
  \label{fig:gw_prob}
\end{figure}

\subsection{Agreement with the Weighted Overlay}
\label{sec:agreement}

Compared pixel by pixel against the weighted-overlay reference labels across all
37~million valid pixels, the XGBoost map agrees on 99.05\% of them
(Figure~\ref{fig:overlay_compare}). This high agreement is the expected
consequence of training on those labels, and it should be interpreted as evidence
that the model faithfully learned the overlay---not as an independent measure of
correctness. The small disagreements are concentrated, as one would expect, along
the class boundaries where the continuous score crosses a percentile threshold.
Establishing whether the map is correct in a stronger sense requires evidence
external to the labels, which the next section provides.

\begin{figure}[H]
  \centering
  \includegraphics[width=\textwidth]{Fig4_Overlay_vs_ML_Comparison}
  \caption{Side-by-side comparison of the weighted-overlay reference labels and
           the XGBoost prediction, which agree on 99.05\% of the 37~million
           valid pixels.}
  \label{fig:overlay_compare}
\end{figure}

\section{Spatial-Realism Validation}
\label{sec:validation_results}

The four independent checks defined in Section~\ref{sec:validation_design} were
applied to the predicted map. Their combined results are shown in
Figure~\ref{fig:spatial_realism} and described below.

\begin{figure}[H]
  \centering
  \includegraphics[width=\textwidth]{spatial_realism_validation}
  \caption{Spatial-realism validation summary: (1) known alluvial sites against
           predicted class; (2) recharge-zone discrimination from independent
           rain-gauge stations; (3) cross-study consistency of the
           High-potential fraction; and (4) the district-scale potential
           gradient.}
  \label{fig:spatial_realism}
\end{figure}

\paragraph{Known productive sites.}
All four published groundwater sites in the Rawalpindi alluvial basin fell within
Medium or High potential ground on the predicted map, with High-class
probabilities at or close to one---a 100\% hit rate for locations independently
known to be productive. The number of such sites is small, so this is a
qualitative rather than statistical result, but it is the most direct available
check and the map passes it cleanly.

\paragraph{Recharge-zone discrimination.}
The independent rain-gauge stations of Khan et al.
\cite{khan2023rainfallrunoff}, split into alluvial and non-alluvial groups by
hydrogeological setting, showed a clear separation in predicted potential. The
alluvial stations had a markedly higher predicted High-probability than the
non-alluvial ones, and a one-sided Mann--Whitney $U$ test confirmed the
difference was statistically significant ($p = 0.023$). Because these stations
played no part in training the model, this is genuine external evidence that the
map distinguishes recharge-favourable from recharge-poor ground.

\paragraph{Cross-study consistency.}
The 20\% of the plateau classified as High potential is consistent in magnitude
with independent AHP-based assessments of the same region. Hafeez et al.
\cite{hafeez2025rwh}, mapping water-storage and rainwater-harvesting suitability
across essentially the same districts, identified a minority of the area as most
suitable, in the same broad range. Both approaches agree that high-potential
ground is the exception rather than the rule across this semi-arid upland, which
argues against the prediction being systematically over- or under-stated.

\paragraph{District-scale gradient.}
Ranking the five districts by their predicted fraction of High potential
reproduces the expected hydrogeological order almost exactly. Rawalpindi, the
wettest and most alluvial district, has the largest High fraction at about 61\%,
descending through Attock and Jhelum to Mianwali and Chakwal at only 4--5\%. The
Spearman rank correlation between predicted and expected ordering is
$\rho = 1.00$, a perfect match. This gradient---wet, alluvial north against dry,
hard-rock south---is the single most robust piece of validation evidence,
because it rests on the relative ordering of large district populations rather
than on a handful of points.

\paragraph{Summary.}
Taken together, the four checks point the same way: the map reproduces known
productive ground, separates recharge zones that were never shown to the model,
matches the magnitude of independent regional studies, and follows the expected
district gradient perfectly. None of these tests depends on the weighted-overlay
labels, so they address the credibility question that internal accuracy cannot.
Their interpretation, and the limits of what they can establish, is the subject
of Chapter~\ref{ch:discussion}.
